% !TeX root = ../main.tex

\documentclass[../main.tex]{subfiles}
\begin{document}

\chapter{TỔNG QUAN}

\section{Thực trạng quản lý CTRSH tại thành phố Hội An}

Thành phố Hội An là đô thị di sản và là trung tâm du lịch quan trọng của tỉnh Quảng Nam. Quá trình phát triển du lịch, dịch vụ lưu trú, nhà hàng, thương mại và gia tăng dân số cơ học làm lượng CTRSH phát sinh có xu hướng tăng nhanh. Theo hồ sơ điều chỉnh công nghệ của dự án, lượng rác thu gom tại Hội An đã tăng từ khoảng 65,5 tấn/ngày năm 2013 lên trên 100 tấn/ngày năm 2019 và được dự báo có thể đạt 120--140 tấn/ngày trong giai đoạn 2025--2030 \cite{BaocaoHoiAn2024}. Với đặc thù là đô thị du lịch, chất thải thực phẩm, bao bì nhựa, giấy, túi nylon và rác phát sinh theo mùa du lịch chiếm tỷ lệ đáng kể.

Kết quả phân loại CTRSH tại Hội An cho thấy rác hữu cơ chiếm trung bình 62,56\% khối lượng, nhựa và bao nylon chiếm 10,16\%, giấy chiếm 7,42\%, đất cát chiếm 7,76\%, thủy tinh chiếm 2,28\% và kim loại chiếm 1,45\%. Về tính chất cháy, nhóm chất thải cháy được chiếm khoảng 86,62\%, còn nhóm không cháy chiếm khoảng 13,38\%. Độ ẩm trung bình của mẫu rác là 56,3\%, độ tro khoảng 8,5\%. Các số liệu này cho thấy rác Hội An có tiềm năng thu hồi năng lượng nhưng bắt buộc phải có tiền xử lý tách tạp chất, giảm ẩm và ổn định kích thước trước khi đưa vào hệ thống xử lý nhiệt.

Nhà máy xử lý rác thải sinh hoạt hiện hữu tại khu vực xã Cẩm Hà sử dụng công nghệ đốt trực tiếp với công suất thiết kế 100 tấn/ngày đêm, nhưng quá trình vận hành không ổn định, công suất thực tế chỉ đạt khoảng 25--30\% công suất thiết kế và chưa đáp ứng yêu cầu nghiệm thu môi trường. Khi hệ thống xử lý tại chỗ không vận hành ổn định, rác phát sinh hằng ngày phải vận chuyển đến bãi rác Tam Nghĩa, huyện Núi Thành, cách Hội An khoảng 100 km. Phương án vận chuyển xa này làm tăng chi phí xử lý, tiêu hao nhiên liệu vận tải, gia tăng nguy cơ rơi vãi, phát tán mùi trên tuyến đường và không bền vững trong dài hạn.

\begin{table}[htbp]
\centering
\caption{Thành phần trung bình CTRSH tại thành phố Hội An theo khối lượng}
\label{tab:thanh-phan-rac-hoi-an}
\setlength{\tabcolsep}{5pt}
% \resizebox{\textwidth}{!}{%
\begin{tabular}{>{\centering\arraybackslash}p{0.08\textwidth}>{\raggedright\arraybackslash}p{0.52\textwidth}>{\centering\arraybackslash}p{0.25\textwidth}}
\toprule
\textbf{TT} & \textbf{Thành phần} & \textbf{Giá trị bình quân, \%} \\
\midrule
1 & Rác hữu cơ & 62,56 \\
2 & Giấy & 7,42 \\
3 & Vải & 1,54 \\
4 & Gỗ, củi, cành cây & 4,72 \\
5 & Nhựa, bao nylon & 10,16 \\
6 & Da và cao su & 0,22 \\
7 & Kim loại & 1,45 \\
8 & Thủy tinh & 2,28 \\
9 & Sành sứ, đất cát, xỉ than, bùn và thành phần khác & 9,65 \\
\midrule
\multicolumn{2}{l}{\textbf{Tổng cộng}} & \textbf{100,00} \\
\bottomrule
\end{tabular}%
% }
\end{table}

% \begin{figure}[htbp]
% \centering
% \fbox{\parbox{0.86\textwidth}{\centering\vspace{2.6cm}
% \textit{bản vẽ hiện trạng sử dụng đất hoặc vị trí nhà máy xử lý CTRSH thành phố Hội An.}
% \vspace{2.6cm}}}
% \caption{Vị trí và hiện trạng khu vực dự án Nhà máy xử lý CTRSH thành phố Hội An}
% \label{fig:vi-tri-du-an}
% \end{figure}

\section{Tổng quan các phương pháp xử lý CTRSH}

Các phương pháp xử lý CTRSH có thể được chia thành ba nhóm lớn gồm xử lý cơ học -- tái chế, xử lý sinh học và xử lý nhiệt. Việc lựa chọn công nghệ cho một địa phương không chỉ phụ thuộc vào công suất xử lý mà còn phụ thuộc trực tiếp vào thành phần rác, độ ẩm, khả năng phân loại tại nguồn, quỹ đất, yêu cầu môi trường và năng lực vận hành thực tế.

Chôn lấp hợp vệ sinh có ưu điểm là tiếp nhận được nhiều loại chất thải và đầu tư ban đầu tương đối thấp, nhưng không phù hợp với định hướng giảm chôn lấp vì tiêu tốn quỹ đất và phát sinh nước rỉ rác, khí bãi rác. Ủ compost tận dụng được thành phần hữu cơ cao trong rác nhưng yêu cầu phân loại tốt, thời gian xử lý dài và phụ thuộc thị trường tiêu thụ sản phẩm phân hữu cơ. Đốt trực tiếp có khả năng giảm nhanh khối lượng và thể tích rác, song đối với rác có độ ẩm cao thì cần bổ sung nhiên liệu hoặc sấy sơ bộ; đồng thời hệ thống xử lý khí thải phải được đầu tư nghiêm ngặt để kiểm soát bụi, kim loại nặng, dioxin và furan \cite{Tait2020}.

Khí hóa là quá trình chuyển hóa vật liệu chứa carbon trong điều kiện thiếu oxy ở nhiệt độ cao, thường trong khoảng 800--1.200$^\circ$C, tạo ra khí tổng hợp gồm chủ yếu \ce{CO}, \ce{H2}, \ce{CH4} và một phần \ce{CO2} \cite{Arena2012}. Về mặt lý thuyết, khí tổng hợp có thể dùng cho động cơ khí, tuabin khí hoặc lò hơi. Tuy nhiên, khí hóa rác sinh hoạt đòi hỏi nguyên liệu RDF có độ ẩm thấp, kích thước và thành phần đồng đều. Dòng khí tổng hợp thường chứa bụi, hắc ín, hơi nước, hợp chất clo, lưu huỳnh và các cấu tử hữu cơ dễ ngưng tụ nên hệ thống làm sạch khí rất phức tạp, đặc biệt khi cấp cho tuabin khí.

Nhiệt phân là quá trình phân hủy nhiệt chất hữu cơ trong điều kiện không có hoặc gần như không có oxy, thường ở dải 350--600$^\circ$C đối với rác nhựa và RDF \cite{Bridgwater2012,Sharuddin2016}. Sản phẩm của quá trình gồm dầu nhiệt phân, khí không ngưng và than carbon. So với khí hóa, nhiệt phân vận hành ở nhiệt độ thấp hơn, không đốt trực tiếp rác trong buồng phản ứng, giảm phát sinh bụi mịn và tạo ra sản phẩm lỏng có thể lưu trữ, chưng cất, pha trộn và sử dụng cho động cơ diesel sau xử lý thích hợp \cite{Miandad2016,Williams2007}. Vì vậy, công nghệ nhiệt phân có mức độ phù hợp cao hơn đối với cấu hình xử lý rác nhựa, nylon và phần RDF đã được tiền xử lý.

\begin{table}[htbp]
\centering
\caption{So sánh khái quát một số phương pháp xử lý CTRSH}
\label{tab:so-sanh-phuong-phap-xu-ly}
\setlength{\tabcolsep}{5pt}
\resizebox{\textwidth}{!}{%
\begin{tabular}{>{\raggedright\arraybackslash}p{0.18\textwidth}p{0.26\textwidth}p{0.32\textwidth}>{\raggedright\arraybackslash}p{0.20\textwidth}}
\toprule
\textbf{Phương pháp} & \textbf{Ưu điểm} & \textbf{Hạn chế} & \textbf{Sản phẩm chính} \\
\midrule
Chôn lấp hợp vệ sinh & Đơn giản, tiếp nhận nhiều loại rác & Chiếm đất, phát sinh nước rỉ rác và khí nhà kính & Khí bãi rác, nước rỉ rác \\
Ủ compost & Phù hợp rác hữu cơ đã phân loại & Thị trường sản phẩm hạn chế, không xử lý rác nhựa & Phân hữu cơ \\
Đốt trực tiếp & Giảm nhanh thể tích rác, có thể thu hồi nhiệt & Yêu cầu kiểm soát khí thải nghiêm ngặt, bất lợi với rác ẩm & Nhiệt, tro xỉ \\
Khí hóa & Thu hồi khí nhiên liệu, hiệu suất lý thuyết cao & Yêu cầu RDF đồng nhất, làm sạch khí phức tạp & Khí tổng hợp, tro xỉ \\
Nhiệt phân & Tạo dầu lỏng dễ lưu trữ, giảm bụi do không đốt trực tiếp & Cần phân loại và tinh chế sản phẩm dầu & Dầu nhiệt phân, khí, than carbon \\
\bottomrule
\end{tabular}%
}
\end{table}

\section{Giới thiệu dự án và công nghệ ban đầu}

Dự án Nhà máy xử lý CTRSH thành phố Hội An có công suất xử lý 120 tấn/ngày đêm, được đầu tư tại thôn Bàu Ốc, xã Cẩm Hà, thành phố Hội An, tỉnh Quảng Nam. Diện tích mặt bằng xây dựng công trình là 14.068,5 m\textsuperscript{2}. Dự án được triển khai nhằm xử lý rác sinh hoạt phát sinh trên địa bàn thành phố, giảm phụ thuộc vào vận chuyển rác đi xa, giảm áp lực chôn lấp và từng bước thu hồi năng lượng từ chất thải.

Theo phương án được phê duyệt ban đầu, dây chuyền công nghệ gồm ba khối chính: sản xuất viên nén RDF từ CTRSH, khí hóa nhiên liệu RDF và phát điện bằng tuabin khí. Sau khi rác được tiếp nhận, xé bao, phân loại, cắt, xử lý nhiệt, ép tách nước và sấy, phần cháy được được tạo thành viên RDF để cấp cho lò khí hóa. Khí tổng hợp sau khi làm sạch được cấp cho cụm tuabin khí -- máy phát điện. Công suất phát điện theo thiết kế ban đầu khoảng 2,6 MW \cite{BaocaoHoiAn2024}.

Trong quá trình rà soát triển khai thực tế, phương án khí hóa -- tuabin khí bộc lộ nhiều điểm không phù hợp với đặc tính rác Hội An. Nhiên liệu RDF đầu vào khó duy trì đồng nhất do rác chưa được phân loại triệt để tại nguồn, độ ẩm cao và thành phần biến động theo mùa. Dòng khí tổng hợp từ rác có nguy cơ chứa hắc ín, bụi, hợp chất clo, lưu huỳnh và hơi dầu nặng; nếu không làm sạch triệt để có thể gây mài mòn, bám bẩn và giảm tuổi thọ tuabin. Ngoài ra, tuabin khí thường tối ưu ở quy mô công suất lớn hơn, trong khi dự án chỉ ở mức vài MW, làm giảm hiệu quả kinh tế của cấu hình thiết bị.

\begin{figure}[H]
    \centering
    \includegraphics[width=0.7\textwidth]{Hinhve/so-do-ban-dau.png}
    \caption{Sơ đồ khái quát phương án công nghệ khí hóa -- tuabin khí ban đầu}
    \label{fig:so-do-cong-nghe-ban-dau}
\end{figure}

\section{Mục tiêu và phạm vi nghiên cứu của đồ án}

Trên cơ sở các vấn đề thực tiễn nêu trên, đồ án tập trung nghiên cứu phương án điều chỉnh công nghệ từ khí hóa -- tuabin khí sang nhiệt phân -- chưng cất dầu -- phát điện bằng động cơ diesel. Mục tiêu kỹ thuật của phương án điều chỉnh là nâng cao tính ổn định vận hành, giảm yêu cầu khắt khe đối với chất lượng khí nhiên liệu, tận dụng sản phẩm dầu nhiệt phân có khả năng lưu trữ, đồng thời giảm khối lượng chất thải phải chôn lấp.

Các nội dung chính của đồ án gồm: phân tích cơ sở lựa chọn công nghệ nhiệt phân đối với điều kiện rác sinh hoạt Hội An; thuyết minh quy trình tiền xử lý, tạo RDF, nhiệt phân, chưng cất dầu, phát điện và xử lý môi trường; tính toán cân bằng vật chất và năng lượng ở mức thiết kế sơ bộ; lựa chọn các nhóm thiết bị chính; đánh giá các yêu cầu quản lý vận hành, an toàn, môi trường và hiệu quả kinh tế -- xã hội. Các số liệu sử dụng trong đồ án được ưu tiên lấy từ hồ sơ điều chỉnh công nghệ dự án, kết quả phân loại rác, bảng cân bằng vật chất, thông số thiết bị và các tài liệu khoa học đã được trích dẫn trong danh mục tài liệu tham khảo.

\begin{sloppypar}
Phạm vi nghiên cứu của đồ án là hệ thống xử lý CTRSH công suất 120 tấn/ngày đêm tại thành phố Hội An.
\end{sloppypar}

\section{Cơ sở dữ liệu và phương pháp tiếp cận nghiên cứu}

Nguồn số liệu chính của đồ án là hồ sơ điều chỉnh công nghệ của dự án Nhà máy xử lý CTRSH thành phố Hội An, bao gồm bảng phân loại thành phần rác, cân bằng vật chất, cân bằng năng lượng, thông số dầu nhiệt phân, danh mục thiết bị và thuyết minh các hệ thống xử lý môi trường \cite{BaocaoHoiAn2024}. Các tài liệu khoa học trong danh mục tham khảo được sử dụng để giải thích bản chất quá trình khí hóa, nhiệt phân, xử lý sản phẩm dầu và các rủi ro công nghệ liên quan \cite{Arena2012,Bridgwater2012,Sharuddin2016,Miandad2016}.

\begin{table}[H]
\centering
\caption{Các nhóm dữ liệu sử dụng trong đồ án}
\label{tab:nhom-du-lieu-do-an}
\setlength{\tabcolsep}{5pt}
\scalebox{0.9}{
\begin{tabular}{>{\raggedright\arraybackslash}p{0.22\textwidth}p{0.36\textwidth}p{0.36\textwidth}}
\toprule
\textbf{Nhóm dữ liệu} & \textbf{Nội dung sử dụng} & \textbf{Vai trò trong đồ án} \\
\midrule
Dữ liệu hiện trạng và dự án & Công suất 120 tấn/ngày đêm, vị trí xây dựng, diện tích mặt bằng, hiện trạng nhà máy đốt cũ, khoảng cách vận chuyển rác đến bãi xử lý tạm thời & Làm rõ tính cấp thiết, quy mô nghiên cứu và các điều kiện biên khi lựa chọn công nghệ \\
Dữ liệu thành phần CTRSH & Thành phần hữu cơ, giấy, nhựa, vải, gỗ, kim loại, thủy tinh, đất cát, độ ẩm, độ tro và tỷ lệ cháy được/không cháy được & Đánh giá khả năng tạo RDF, yêu cầu tiền xử lý, độ phù hợp của khí hóa và nhiệt phân \\
Dữ liệu cân bằng vật chất -- năng lượng & Khối lượng sau phân loại, bức xạ nhiệt, ép, sấy, lượng RDF cấp lò, tỷ lệ dầu/khí/than, nhiệt trị RDF và hiệu suất chuyển đổi & Làm cơ sở tính toán, xác định dòng sản phẩm và quy mô nhóm thiết bị chính \\
Dữ liệu môi trường và vận hành & Hệ thống xử lý khí thải, nước thải, chất thải rắn, chương trình quan trắc, yêu cầu PCCC, cảm biến LEL, đào tạo vận hành & Làm cơ sở cho phần quản lý vận hành và kiểm soát môi trường trong Chương 4 \\
\bottomrule
\end{tabular}%
}
\end{table}

Phương pháp tiếp cận của đồ án là phân tích, luận chứng và tính toán phương án điều chỉnh công nghệ trên nền dự án đã có, không phải thiết kế mới hoàn toàn tách rời thực tế. Do đó, các thông số đã được xác định trong hồ sơ như công suất 120 tấn/ngày đêm, vị trí tại thôn Bàu Ốc, diện tích mặt bằng 14.068,5 m\textsuperscript{2}, các dòng vật chất sau phân loại, ép, sấy và lượng RDF cấp lò nhiệt phân được lấy làm điều kiện biên tính toán. Cách tiếp cận này phản ánh mục tiêu của thiết kế sơ bộ: đánh giá tính hợp lý kỹ thuật của phương án, tổ chức hệ thống tính toán, làm rõ cơ sở lựa chọn thiết bị và xác định yêu cầu vận hành -- kiểm soát môi trường.
\begin{figure}[H]
\centering
\scalebox{0.9}{
\begin{tikzpicture}[
    node distance=0.6cm and 0.5cm,
    stepnum/.style={
        circle,
        draw=teal!80!black,
        fill=teal!60!black,
        text=white,
        font=\bfseries\Large,
        minimum size=0.9cm,
        inner sep=0pt
    },
    box/.style={
        rectangle,
        draw=teal!70!black,
        fill=teal!5,
        thick,
        rounded corners=6pt,
        text width=10.5cm,
        align=left,
        inner sep=8pt,
        font=\small
    },
    arrow/.style={
        -{Stealth[length=2.5mm]},
        ultra thick,
        draw=teal!60!black
    }
]
\node (b1) [box] {\textbf{\color{teal!80!black}Dữ liệu đầu vào}\\ $\bullet$ Công suất xử lý chất thải rắn sinh hoạt: 120 tấn/ngày đêm.\\ $\bullet$ Đặc tính rác Hội An: Tỷ lệ hữu cơ cao, độ ẩm cao (56,3\%), nhựa (10,16\%).\\ $\bullet$ Ràng buộc mặt bằng (14.068,5 m\textsuperscript{2}) và hạ tầng đấu nối điện, nước.};
\node (s1) [stepnum, left=of b1] {1};

\node (b2) [box, below=of b1] {\textbf{\color{teal!80!black}Phân tích \& Luận chứng công nghệ}\\ $\bullet$ Đánh giá ưu/nhược điểm các phương pháp (chôn lấp, đốt, khí hóa, nhiệt phân).\\ $\bullet$ Luận chứng sự chuyển dịch từ khí hóa -- tuabin khí sang nhiệt phân -- diesel.\\ $\bullet$ Xác định cấu hình công nghệ: Tiền xử lý $\rightarrow$ Tạo RDF $\rightarrow$ Nhiệt phân $\rightarrow$ Diesel.};
\node (s2) [stepnum, left=of b2] {2};

\node (b3) [box, below=of b2] {\textbf{\color{teal!80!black}Tính toán cân bằng vật chất \& năng lượng}\\ $\bullet$ Cân bằng vật chất: Theo dõi các dòng rác, RDF, dầu nhiệt phân, than carbon, tro xỉ.\\ $\bullet$ Cân bằng năng lượng: Xác định nhiệt trị RDF, hiệu suất nhiệt phản ứng, điện năng phát.};
\node (s3) [stepnum, left=of b3] {3};

\node (b4) [box, below=of b3] {\textbf{\color{teal!80!black}Thiết kế sơ bộ, lựa chọn thiết bị \& Đánh giá vận hành, môi trường}\\ $\bullet$ Xác định quy mô, thông số lò quay nhiệt phân, tháp chưng cất và tổ máy phát điện diesel.\\ $\bullet$ Thiết kế hệ thống xử lý khí thải đạt QCVN 19:2024 và nước thải đạt QCVN 40:2025.\\ $\bullet$ Đề xuất giải pháp an toàn công nghiệp, PCCC và quản lý chất thải thứ cấp.};
\node (s4) [stepnum, left=of b4] {4};

\draw [arrow] (s1) -- (s2);
\draw [arrow] (s2) -- (s3);
\draw [arrow] (s3) -- (s4);
\end{tikzpicture}
}
\caption{Trình tự tiếp cận nghiên cứu trong đồ án}
\label{fig:phuong-phap-nghien-cuu}
\end{figure}
Trong quá trình sử dụng số liệu, đồ án phân biệt ba nhóm dữ liệu. Nhóm thứ nhất là dữ liệu thiết kế ràng buộc, chẳng hạn công suất tiếp nhận rác, diện tích mặt bằng, tỷ lệ thành phần rác, độ ẩm và cân bằng khối lượng. Nhóm thứ hai là dữ liệu công nghệ dùng để tính toán sơ bộ, gồm tỷ lệ sản phẩm nhiệt phân, nhiệt trị RDF, hiệu suất chuyển đổi năng lượng và sản lượng dầu sau chưng cất. Nhóm thứ ba là dữ liệu quản lý vận hành và môi trường, bao gồm giới hạn khí thải, nước thải, chương trình giám sát, yêu cầu an toàn cháy nổ và quản lý chất thải thứ cấp. Việc tách nhóm dữ liệu giúp tránh nhầm lẫn giữa thông số đã có cơ sở trong hồ sơ và thông số cần được kiểm chứng trong vận hành thử nghiệm.

Về trình tự nghiên cứu, đồ án đi từ vấn đề thực tiễn đến cơ sở công nghệ, sau đó thực hiện tính toán thiết kế sơ bộ và đánh giá vận hành, bảo vệ môi trường. Trình tự này phản ánh đúng logic của một dự án xử lý chất thải: trước hết phải xác định bài toán rác, tiếp theo xác định công nghệ có khả năng xử lý dòng rác đó, sau đó kiểm tra bằng cân bằng vật chất -- năng lượng và cuối cùng là thiết kế sơ bộ thiết bị kết hợp đánh giá vận hành, môi trường.

\section{Phân tích đặc tính CTRSH Hội An dưới góc nhìn thiết kế công nghệ}

Đối với công nghệ xử lý nhiệt, CTRSH không thể được xem như một loại nhiên liệu đồng nhất. Rác sinh hoạt là hỗn hợp biến động liên tục theo nguồn phát sinh, mùa vụ, thói quen tiêu dùng, mức độ phân loại tại nguồn và phương thức thu gom. Tại Hội An, yếu tố du lịch làm cho dòng rác có đặc trưng phức tạp hơn so với đô thị thuần dân cư. Rác từ nhà hàng, khách sạn, chợ, khu dân cư và hoạt động thương mại có tỷ lệ hữu cơ cao, lẫn nhiều bao bì nhựa, giấy, vải, thủy tinh, kim loại và tạp chất vô cơ. Vì vậy, việc lựa chọn công nghệ không chỉ căn cứ vào giá trị trung bình của thành phần rác, mà còn phải xét đến khả năng chịu dao động của dây chuyền.

Kết quả phân loại trong hồ sơ dự án cho thấy rác hữu cơ là thành phần chiếm tỷ lệ lớn nhất, đạt bình quân 62,56\% khối lượng; giấy chiếm 7,42\%; vải chiếm 1,54\%; gỗ, củi, cành cây chiếm 4,72\%; nhựa và bao nylon chiếm 10,16\%; da và cao su chiếm 0,22\%. Tổng phần cháy được đạt 86,62\%, trong khi phần không cháy được gồm kim loại, thủy tinh, sành sứ, đất cát, xỉ than, bùn và các thành phần khác chiếm 13,38\% \cite{BaocaoHoiAn2024}. Về mặt năng lượng, tỷ lệ cháy được cao là điều kiện thuận lợi để thu hồi năng lượng. Tuy nhiên, về mặt vận hành, tỷ lệ hữu cơ và độ ẩm cao lại làm giảm nhiệt trị, tăng tải sấy và làm tăng nguy cơ biến động chất lượng RDF.

Độ ẩm bình quân của CTRSH Hội An được xác định khoảng 56,3\%, trong khi độ tro khoảng 8,5\%. Độ ẩm cao ảnh hưởng trực tiếp đến quá trình xử lý nhiệt theo ba cơ chế. Thứ nhất, một phần đáng kể năng lượng đầu vào phải dùng để nâng nhiệt và hóa hơi nước, làm giảm năng lượng hữu ích cho phản ứng nhiệt phân hoặc khí hóa. Thứ hai, hơi nước làm thay đổi cân bằng nhiệt và thành phần sản phẩm khí, khiến việc kiểm soát nhiệt độ trong lò trở nên khó khăn hơn. Thứ ba, rác ẩm dễ bám dính, kết khối, lên men và phát sinh mùi trong khu tiếp nhận, gây ảnh hưởng đến thiết bị cấp liệu và môi trường lao động. Vì vậy, khối tiền xử lý gồm phân loại, bức xạ nhiệt, ép tách nước và sấy không phải là phần phụ trợ đơn giản mà là điều kiện quyết định cho khả năng vận hành của toàn dây chuyền.

Tỷ lệ đất cát, thủy tinh, sành sứ và kim loại trong rác tuy không lớn so với thành phần hữu cơ, nhưng có ý nghĩa quan trọng đối với thiết bị. Các thành phần này không tham gia tạo dầu hoặc khí có ích, đồng thời gây mài mòn dao cắt, máy nghiền, vít tải, băng tải và lớp chịu nhiệt của lò. Kim loại có thể gây kẹt cơ khí hoặc tia lửa khi va đập; thủy tinh và sành sứ làm tăng hàm lượng tro; đất cát làm tăng bụi và cặn rắn. Do đó, công đoạn phân loại cơ học, tách từ, sàng và loại bỏ tạp chất trơ cần được xem là một phần của thiết kế công nghệ chứ không chỉ là công đoạn vệ sinh nguyên liệu.

\begin{table}[H]
\centering
\caption{Ý nghĩa kỹ thuật của các đặc tính CTRSH Hội An đối với lựa chọn công nghệ}
\label{tab:y-nghia-dac-tinh-rac}
\setlength{\tabcolsep}{5pt}
\scalebox{0.95}{
\begin{tabular}{>{\raggedright\arraybackslash}p{0.22\textwidth}p{0.36\textwidth}p{0.36\textwidth}}
\toprule
\textbf{Đặc tính rác} & \textbf{Ảnh hưởng kỹ thuật} & \textbf{Hàm ý đối với phương án công nghệ} \\
\midrule
Rác hữu cơ và độ ẩm cao & Giảm nhiệt trị, tăng năng lượng sấy, phát sinh nước thải ép và mùi trong khu tiếp nhận & Cần tiền xử lý ẩm hiệu quả; công nghệ phía sau phải chịu được dao động độ ẩm và không yêu cầu RDF quá đồng nhất \\
Tỷ lệ nhựa, nylon và giấy đáng kể & Tạo tiềm năng sinh dầu và khí khi nhiệt phân; đồng thời có thể phát sinh hợp chất chứa clo nếu lẫn PVC & Nhiệt phân phù hợp để thu hồi dầu, nhưng cần kiểm soát phân loại và xử lý khí axit trong hệ thống khí thải \\
Thành phần không cháy khoảng 13,38\% & Làm tăng tro, cặn, bụi, mài mòn và nguy cơ kẹt thiết bị & Cần tách kim loại, thủy tinh, đất cát trước khi tạo RDF và trước khi cấp vào lò nhiệt phân \\
Biến động theo nguồn và mùa & Làm thay đổi nhiệt trị, độ ẩm, tỷ lệ dầu/khí/than và tải xử lý môi trường & Cần bố trí hầm/bunker đệm, trộn đều nguyên liệu và vận hành theo dải thông số thay vì một điểm cố định \\
Chưa phân loại triệt để tại nguồn & Làm tăng tạp chất nguy hại, rác trơ và rủi ro môi trường thứ cấp & Dây chuyền phải có công đoạn phân loại cơ học, thu gom chất thải nguy hại lẫn trong rác và quản lý chất thải thứ cấp \\
\bottomrule
\end{tabular}%
}
\end{table}

Từ góc nhìn thiết kế, thành phần cháy được cao không đồng nghĩa với việc có thể đưa trực tiếp rác vào thiết bị nhiệt. Nếu rác chưa được làm khô, đồng nhất kích thước và loại bỏ tạp chất trơ, quá trình nhiệt sẽ tiêu hao nhiều năng lượng, khó ổn định và phát sinh nhiều sự cố. Đặc biệt, khí hóa yêu cầu vùng phản ứng ở nhiệt độ cao, có cấp tác nhân oxy hóa và cần khí sản phẩm tương đối sạch để cấp cho tuabin khí. Với rác ẩm và không đồng nhất, việc giữ ổn định vùng phản ứng khí hóa và kiểm soát hắc ín là rất khó. Ngược lại, nhiệt phân có thể làm việc ở nhiệt độ thấp hơn, sản phẩm năng lượng chính là dầu lỏng có thể lưu chứa, và khí không ngưng có thể hồi lưu làm nhiên liệu gia nhiệt. Đây là điểm khác biệt có ý nghĩa lớn khi xử lý dòng CTRSH chưa được phân loại triệt để.

Một khía cạnh cần nhấn mạnh là công nghệ nhiệt phân không làm mất yêu cầu tiền xử lý. Nhiệt phân chỉ phù hợp hơn với rác có dao động nhất định, nhưng vẫn cần RDF đủ khô, kích thước phù hợp và tạp chất trơ được giảm đến mức chấp nhận được. Nếu cấp rác quá ẩm vào lò, lượng hơi nước sinh ra sẽ làm giảm nhiệt độ vùng phản ứng, kéo dài thời gian lưu cần thiết và làm giảm chất lượng dầu. Nếu nguyên liệu quá nhiều đất cát, chất rắn sau nhiệt phân tăng lên và làm giảm hiệu suất thu hồi năng lượng. Vì vậy, luận chứng lựa chọn nhiệt phân phải đi kèm với luận chứng về hệ thống RDF.

\section{Ràng buộc thiết kế và yêu cầu kỹ thuật phương án}

\subsection{Ràng buộc quy mô, mặt bằng và hạ tầng}

Công suất 120 tấn/ngày đêm đặt ra yêu cầu nhà máy phải có khả năng tiếp nhận, phân loại và xử lý liên tục trong điều kiện rác phát sinh hằng ngày. Ở quy mô này, nếu giả thiết vận hành 24 giờ/ngày, lưu lượng trung bình của rác đầu vào là 5 tấn/giờ; tuy nhiên trong thực tế, rác được vận chuyển theo ca, theo tuyến thu gom và theo thời điểm tập kết, nên lưu lượng tức thời tại khu tiếp nhận có thể cao hơn nhiều so với giá trị trung bình. Do đó, khu tiếp nhận, hầm chứa, băng tải và thiết bị phân loại phải được bố trí với khả năng đệm nhất định để tránh tình trạng quá tải cục bộ.

% Diện tích mặt bằng 14.068,5 m\textsuperscript{2} là một điều kiện biên quan trọng. Với diện tích này, nhà máy phải bố trí đồng thời khu tiếp nhận rác, tuyến phân loại, khu sản xuất RDF, khu sấy, lò nhiệt phân, hệ chưng cất -- tinh chế dầu, bồn chứa dầu, khu phát điện diesel, hệ xử lý khí thải, hệ xử lý nước thải, kho chất thải rắn, kho chất thải nguy hại, nhà điều hành, đường giao thông nội bộ và khu an toàn PCCC. Mặt bằng không quá lớn khiến việc bố trí phải ưu tiên dòng vật chất một chiều, giảm giao cắt giữa xe rác, sản phẩm dầu, chất thải thứ cấp và tuyến bảo trì thiết bị.

Trong phương án công nghệ ban đầu, hệ khí hóa -- tuabin khí yêu cầu không gian cho lò khí hóa, hệ làm sạch khí, thiết bị lọc bụi, làm mát, xử lý hắc ín, máy nén hoặc thiết bị điều áp nếu cần, tuabin khí và hệ phụ trợ. Khi điều chỉnh sang nhiệt phân -- dầu -- diesel, một phần diện tích được chuyển cho hệ ngưng tụ, chưng cất, tinh chế và bồn chứa dầu. Sự thay đổi này không chỉ là thay đổi thiết bị mà còn làm thay đổi triết lý bố trí mặt bằng: dòng năng lượng không còn truyền trực tiếp từ khí nóng sang tuabin, mà đi qua sản phẩm trung gian là dầu lỏng. Vì dầu có thể lưu chứa, mặt bằng cần có khu bồn, đê bao chống tràn, khoảng cách an toàn và tuyến đường phục vụ PCCC.

\begin{figure}[htbp]
\centering
\fbox{\parbox{0.88\textwidth}{\centering\vspace{2.8cm}
\textit{sơ đồ dòng vật chất trên mặt bằng: rác đầu vào -- phân loại -- RDF -- nhiệt phân -- dầu -- phát điện -- chất thải thứ cấp.}
\vspace{2.8cm}}}
\caption{Quan hệ giữa dòng vật chất và bố trí mặt bằng nhà máy}
\label{fig:dong-vat-chat-mat-bang}
\end{figure}

Ràng buộc hạ tầng còn thể hiện ở khả năng đấu nối điện, cấp nước, thoát nước và giao thông. Hệ thống phát điện diesel sử dụng dầu nhiệt phân có thể cấp điện tự dùng cho nhà máy, nhưng vẫn cần đấu nối với lưới để khởi động, dự phòng và đảm bảo vận hành trong trường hợp dây chuyền nhiệt phân chưa phát đủ dầu. Hệ thống cấp nước chủ yếu phục vụ vệ sinh, làm mát tuần hoàn, xử lý khí thải và PCCC. Nước thải sản xuất phát sinh đáng kể ở công đoạn ép tách nước và vệ sinh khu tiếp nhận, nên tuyến thu gom nước thải phải được tách khỏi nước mưa sạch để giảm tải cho trạm xử lý.

Đối với giao thông nội bộ, yêu cầu quan trọng là xe rác vào nhà máy không đi chung tuyến với xe vận chuyển dầu, chất thải nguy hại hoặc vật tư hóa chất nếu có thể tránh được. Dòng xe rác thường mang theo nước rỉ, mùi và tạp chất bám ngoài, trong khi khu dầu và phát điện yêu cầu kiểm soát nguồn lửa, tĩnh điện và vệ sinh công nghiệp. Vì vậy, tổ chức mặt bằng cần hỗ trợ quản lý rủi ro, không chỉ hỗ trợ vận chuyển. Các khu vực có nguy cơ cao như bồn dầu, lò nhiệt phân, phòng máy phát và kho hóa chất phải có lối tiếp cận cho lực lượng PCCC và có khả năng cô lập khi xảy ra sự cố.

\subsection{Yêu cầu kỹ thuật đối với phương án điều chỉnh}

Phương án điều chỉnh từ khí hóa -- tuabin khí sang nhiệt phân -- chưng cất dầu -- động cơ diesel cần thỏa mãn đồng thời bốn nhóm yêu cầu: xử lý được rác địa phương, vận hành ổn định, kiểm soát môi trường và có hiệu quả sử dụng năng lượng. Nếu chỉ xét riêng hiệu suất lý thuyết, nhiều công nghệ xử lý nhiệt có thể cho kết quả hấp dẫn; nhưng trong điều kiện rác ẩm, chưa phân loại tốt và quy mô công suất vừa, yêu cầu ổn định vận hành thường quan trọng hơn hiệu suất danh nghĩa.

\begin{table}[H]
\centering
\caption{Các yêu cầu kỹ thuật chính đối với phương án nhiệt phân -- dầu -- diesel}
\label{tab:yeu-cau-phuong-an-dieu-chinh}
\setlength{\tabcolsep}{5pt}
\scalebox{0.95}{
\begin{tabular}{>{\raggedright\arraybackslash}p{0.22\textwidth}p{0.36\textwidth}p{0.36\textwidth}}
\toprule
\textbf{Nhóm yêu cầu} & \textbf{Nội dung kỹ thuật} & \textbf{Ý nghĩa đối với dự án} \\
\midrule
Nguyên liệu RDF & Giảm ẩm, giảm kích thước, tách tạp chất trơ, duy trì cấp liệu đều & Ổn định nhiệt độ lò, giảm kẹt thiết bị và nâng cao tỷ lệ thu hồi dầu \\
Lò nhiệt phân & Làm việc trong điều kiện thiếu oxy, kiểm soát nhiệt độ 350--600\textdegree C, đảo trộn tốt và thời gian lưu phù hợp & Chuyển hóa phần hữu cơ và nhựa thành dầu, khí và than carbon với mức dao động chấp nhận được \\
Ngưng tụ và chưng cất & Thu hồi hơi dầu, tách dầu nhẹ/dầu nặng, giảm nước và cặn & Tạo nhiên liệu lỏng có khả năng lưu chứa và sử dụng cho máy phát diesel sau tinh chế \\
Phát điện diesel & Cấp nhiên liệu ổn định, lọc dầu, kiểm soát khí thải động cơ, dự phòng khởi động bằng nguồn ngoài & Tận dụng năng lượng tại chỗ và giảm phụ thuộc vào điện lưới trong vận hành nhà máy \\
An toàn và môi trường & Cảm biến LEL, thu gom khí kín, xử lý khí thải, xử lý nước thải, PCCC khu dầu, quản lý chất thải thứ cấp & Giảm rủi ro cháy nổ, phát tán mùi, ô nhiễm nước và vi phạm quy chuẩn môi trường \\
\bottomrule
\end{tabular}%
}
\end{table}

Yêu cầu đầu tiên là khả năng thích nghi với RDF có chất lượng dao động. Hệ thống phải cho phép phối trộn, lưu chứa đệm và điều chỉnh tốc độ cấp liệu theo độ ẩm, kích thước và nhiệt trị thực tế. Trong lò nhiệt phân, tốc độ quay, thời gian lưu, nhiệt độ buồng phản ứng và chế độ gia nhiệt cần được điều chỉnh trong một dải nhất định để tránh tình trạng nguyên liệu chưa phân hủy hết hoặc bị quá nhiệt cục bộ. Khả năng thích nghi này đặc biệt quan trọng trong giai đoạn mùa mưa, khi độ ẩm rác tăng và tỷ lệ hữu cơ ướt trong rác đầu vào cao hơn.

Yêu cầu thứ hai là kiểm soát kín dòng khí và hơi dầu. Nhiệt phân tạo ra hơi dầu, khí không ngưng và một phần hợp chất hữu cơ bay hơi. Nếu hệ thống ngưng tụ, đường ống, bồn chứa và tháp xử lý không kín, nguy cơ phát tán mùi, VOC và khí dễ cháy sẽ tăng. Vì vậy, phương án công nghệ phải coi hệ thống thu gom hơi dầu và khí không ngưng là tuyến an toàn chính. Khí không ngưng cần được hồi lưu về buồng đốt hoặc xử lý triệt để, không được xả trực tiếp ra môi trường.

Yêu cầu thứ ba là tạo ra sản phẩm dầu có thể sử dụng cho động cơ diesel sau xử lý thích hợp. Dầu nhiệt phân thô thường chứa nước, cặn rắn, hợp chất oxy hóa, hợp chất không no, lưu huỳnh, clo và phân đoạn nặng. Nếu đưa trực tiếp vào động cơ, dầu có thể gây ăn mòn, đóng cặn kim phun, tăng khói, làm giảm tuổi thọ bơm cao áp và buồng đốt. Do đó, hệ chưng cất và tinh chế dầu phải có vai trò nâng cấp nhiên liệu, tách phần nhẹ, kiểm soát nước/cặn và tuần hoàn dầu nặng về lò hoặc xử lý theo quy định.

Yêu cầu thứ tư là kiểm soát môi trường theo hướng chủ động. Nhà máy không chỉ cần xử lý khí thải ở cuối đường ống, mà phải giảm phát sinh ô nhiễm ngay từ nguồn. Điều này bao gồm kiểm soát độ ẩm RDF để giảm hơi nước và khí mùi, loại bỏ tạp chất chứa clo nếu có thể, duy trì buồng đốt khí không ngưng ở nhiệt độ phù hợp, làm nguội nhanh khí thải để hạn chế tái hình thành dioxin/furan, thu gom nước thải ép rác, quản lý cặn dầu và chất thải nguy hại.

% \section{Cấu trúc đồ án}

% Đồ án triển khai theo trình tự: Chương~2 trình bày cơ sở lý thuyết và luận chứng lựa chọn công nghệ; Chương~3 tập trung tính toán cân bằng vật chất -- năng lượng và thiết kế sơ bộ thiết bị; Chương~4 đánh giá vận hành, môi trường và hiệu quả dự án. Quan điểm xuyên suốt là ưu tiên tính ổn định vận hành, coi tiền xử lý RDF là điều kiện của thu hồi năng lượng, và xem dầu nhiệt phân là nhiên liệu kỹ thuật phải đạt chỉ tiêu trước khi cấp cho động cơ diesel.

% \section{Kết luận chương}

% Chương 1 đã làm rõ bối cảnh và điều kiện đầu vào của đồ án. CTRSH Hội An có hữu cơ cao (62,56\%), độ ẩm lớn (56,3\%) và thành phần biến động, trong khi nhà máy hiện hữu vận hành không ổn định và phải vận chuyển rác xa. Phương án khí hóa -- tuabin khí ban đầu không phù hợp với chất lượng RDF thực tế và quy mô phát điện vài MW. Đồ án chọn hướng nhiệt phân -- chưng cất dầu -- phát điện diesel ở công suất 120 tấn/ngày đêm, trên mặt bằng 14.068,5 m\textsuperscript{2} tại xã Cẩm Hà. Các số liệu thành phần rác, cân bằng khối lượng và thông số thiết bị từ hồ sơ dự án là cơ sở cho tính toán kỹ thuật ở các chương tiếp theo.

% \section{Tóm tắt chương 1}

% Chương 1 đã trình bày bức tranh tổng thể về hiện trạng quản lý CTRSH tại thành phố Hội An và cơ sở lựa chọn phương án điều chỉnh công nghệ. Các nội dung chính bao gồm:

% \begin{itemize}
%     \item \textbf{Hiện trạng Hội An:} Lượng CTRSH Hội An đạt trên 100 tấn/ngày và tiếp tục tăng do phát triển du lịch. Nhà máy xử lý hiện hữu dùng công nghệ đốt trực tiếp vận hành không ổn định, thực tế chỉ đạt 25--30\% công suất thiết kế và chưa đáp ứng yêu cầu môi trường, dẫn đến phải vận chuyển rác xa tới bãi Tam Nghĩa.

%     \item \textbf{Lựa chọn công nghệ:} Sau phân tích so sánh, phương án nhiệt phân -- dầu -- động cơ diesel được luận chứng phù hợp hơn so với phương án khí hóa -- tuabin khí ban đầu, nhờ tính ổn định vận hành cao hơn với rác ẩm và biến động, sản phẩm dầu có thể lưu trữ và kiểm soát chất lượng, mô-đun hóa phù hợp và dễ bảo trì hơn trong điều kiện địa phương.

%     \item \textbf{Phạm vi và quan điểm:} Đồ án tiếp cận theo hướng phân tích phương án điều chỉnh công nghệ trên nền dự án thực tế, ưu tiên tính ổn định vận hành, coi tiền xử lý là một phần không tách rời của công nghệ thu hồi năng lượng, đánh giá dầu nhiệt phân như nhiên liệu kỹ thuật và kiểm soát môi trường gắn với kiểm soát quá trình.
% \end{itemize}

% Các nội dung và số liệu phân tích trong Chương 1 là nền tảng dữ liệu đầu vào và cơ sở luận chứng cho toàn bộ các chương tiếp theo của đồ án.

% \section{Phân tích điều kiện kỹ thuật và phi kỹ thuật cho sự chuyển dịch phương án}

% \subsection{Điều kiện kỹ thuật quyết định sự chuyển dịch}

% Việc điều chỉnh từ phương án khí hóa -- tuabin khí sang phương án nhiệt phân -- dầu -- động cơ diesel không chỉ dựa trên cơ sở so sánh lý thuyết mà còn xuất phát từ các điều kiện thực tế đã được xác nhận qua quá trình rà soát thiết kế và tham khảo dự án tương tự. Có ba nhóm điều kiện kỹ thuật quan trọng dẫn đến quyết định này.

% Nhóm điều kiện thứ nhất liên quan đến chất lượng nguyên liệu RDF. Khí hóa đòi hỏi RDF với độ ẩm không vượt quá 20\%, nhiệt trị ổn định trong khoảng 14--18 MJ/kg và thành phần chlorine dưới 0,5\% để tránh ăn mòn vật liệu chịu nhiệt. Trong thực tế vận hành tại Hội An, rác đầu vào có độ ẩm thường xuyên vượt 50\%, và hệ thống tiền xử lý khó duy trì RDF đầu ra dưới 20\% ẩm trong mùa mưa. Nhiệt phân chấp nhận được RDF có độ ẩm cao hơn (đến 15--20\%) và không đòi hỏi cấp tác nhân oxy hóa, do đó ít nhạy cảm hơn với biến động chất lượng nguyên liệu \cite{Sharuddin2016}.

% Nhóm điều kiện thứ hai liên quan đến hệ thống làm sạch khí sản phẩm. Tuabin khí thương mại đặt yêu cầu rất khắt khe về chất lượng khí nhiên liệu: hàm lượng bụi nhỏ hơn 5 mg/Nm\textsuperscript{3}, hắc ín (tar) dưới 0,05 mg/Nm\textsuperscript{3}, và các cấu tử ăn mòn như \ce{HCl}, \ce{H2S} phải dưới ngưỡng ppm. Đạt được các ngưỡng này từ khí tổng hợp của rác hỗn hợp đòi hỏi hệ thống làm sạch nhiều bậc với chi phí đầu tư lớn và khó vận hành ổn định. Với nhiệt phân, sản phẩm năng lượng là dầu lỏng -- có thể lọc, lắng và tinh chế đến chất lượng phù hợp với động cơ diesel thông thường bằng công nghệ hóa lý quen thuộc \cite{Miandad2016}.


% Nhóm điều kiện thứ ba liên quan đến quy mô và hiệu quả kinh tế. Tuabin khí thương mại đạt hiệu suất kinh tế cao nhất ở quy mô phát điện từ 10 MW trở lên. Tại mức công suất 2--3 MW như của dự án Hội An, tỷ lệ chi phí bảo trì, phụ tùng và nhân lực vận hành cao áp trên mỗi kWh điện phát ra sẽ rất lớn. Ngược lại, động cơ diesel đốt nội có công suất mô-đun 200--500 kW phổ biến trên thị trường, phụ tùng dễ kiếm và đội ngũ bảo trì sẵn có tại địa phương \cite{BaocaoHoiAn2024}.

% \subsection{Điều kiện phi kỹ thuật hỗ trợ quyết định điều chỉnh}

% Bên cạnh các yếu tố kỹ thuật, một số điều kiện phi kỹ thuật cũng đóng vai trò quan trọng trong quyết định điều chỉnh phương án công nghệ.

% \textbf{Năng lực nhân lực vận hành:} Vận hành tuabin khí đòi hỏi kỹ sư có chuyên môn cao về động cơ nhiệt khí, kiểm soát cháy và tự động hóa. Tại địa phương có quy mô vừa như Hội An, việc tuyển dụng và duy trì đội ngũ vận hành như vậy sẽ gặp khó khăn. Trong khi đó, vận hành động cơ diesel là kỹ năng phổ biến trong ngành cơ khí và điện tại Việt Nam, giảm đáng kể yêu cầu về nhân lực chuyên môn cao.

% \textbf{Chuỗi cung ứng phụ tùng:} Phụ tùng tuabin khí chuyên dụng thường phải nhập khẩu với thời gian chờ dài (3--6 tháng), trong khi phụ tùng động cơ diesel có thể mua trong nước hoặc từ các nhà cung cấp khu vực với thời gian chờ ngắn (1--4 tuần). Điều này ảnh hưởng trực tiếp đến khả năng duy trì nhà máy hoạt động liên tục và chi phí tồn kho phụ tùng.

% \textbf{Khả năng huy động vốn và phân kỳ đầu tư:} Hệ thống nhiệt phân -- dầu -- diesel có thể được đầu tư theo từng mô-đun (mỗi lò nhiệt phân và máy phát được lắp đặt riêng biệt), cho phép phân kỳ vốn đầu tư và giảm rủi ro tài chính. Cấu hình khí hóa -- tuabin khí thường phải đầu tư tổng thể ngay từ đầu do tính phụ thuộc cao giữa các thiết bị.

% \subsection{Đánh giá tổng hợp điều kiện áp dụng và khuyến nghị}

% Từ phân tích trên, có thể kết luận rằng phương án nhiệt phân -- dầu -- động cơ diesel phù hợp với điều kiện của dự án Hội An hơn phương án ban đầu theo sáu tiêu chí chính:

% \begin{enumerate}
%     \item Thích nghi với chất lượng nguyên liệu RDF thực tế tại địa phương (ẩm cao, biến động theo mùa).
%     \item Hệ thống xử lý sản phẩm năng lượng (tinh chế dầu) đơn giản hơn so với làm sạch khí tổng hợp.
%     \item Sản phẩm dầu lỏng có thể lưu trữ, kiểm tra chất lượng và sử dụng linh hoạt.
%     \item Phù hợp với quy mô công suất phát điện 2--3 MW (mô-đun diesel phổ biến).
%     \item Yêu cầu nhân lực và phụ tùng phù hợp với năng lực địa phương.
%     \item Cho phép phân kỳ đầu tư theo từng mô-đun nhiệt phân và phát điện.
% \end{enumerate}

% \begin{table}[htbp]
% \centering
% \caption{Tổng hợp so sánh phương án khí hóa -- tuabin khí và phương án nhiệt phân -- dầu -- diesel}
% \label{tab:so-sanh-toan-dien-hai-phuong-an}
% \setlength{\tabcolsep}{5pt}
% \resizebox{\textwidth}{!}{%
% \begin{tabular}{>{\raggedright\arraybackslash}p{0.22\textwidth}>{\centering\arraybackslash}p{0.10\textwidth}p{0.32\textwidth}p{0.32\textwidth}}
% \toprule
% \textbf{Tiêu chí đánh giá} & \textbf{Trọng số} & \textbf{Khí hóa -- Tuabin khí} & \textbf{Nhiệt phân -- Dầu -- Diesel} \\
% \midrule
% Chấp nhận biến động nguyên liệu & Cao & Yêu cầu RDF chất lượng cao, ổn định & Chấp nhận biến động nhất định, ít nhạy cảm hơn \\
% Độ phức tạp vận hành và bảo trì & Cao & Rất phức tạp (tuabin khí, làm sạch khí nhiều bậc) & Trung bình (động cơ diesel thông dụng, tinh chế dầu) \\
% Phù hợp quy mô 2--3 MW & Cao & Không tối ưu (tuabin thường lớn hơn) & Phù hợp (nhiều máy phát mô-đun song song) \\
% Khả năng lưu trữ sản phẩm năng lượng & Trung bình & Không lưu trữ được khí (dùng ngay) & Lưu trữ được dầu (bồn chứa) \\
% Chi phí làm sạch sản phẩm & Trung bình & Rất cao (nhiều bậc lọc khí) & Trung bình (chưng cất, tinh chế dầu) \\
% Năng lực nhân lực địa phương & Trung bình & Đòi hỏi chuyên môn tuabin khí đặc thù & Phổ biến (cơ khí, điện diesel thông thường) \\
% Phân kỳ đầu tư & Thấp & Khó phân kỳ (hệ thống tổng thể) & Dễ phân kỳ (từng mô-đun lò + máy phát) \\
% \bottomrule
% \end{tabular}%
% }
% \end{table}

% \section{Dự báo nhu cầu xử lý và năng lực cần thiết trong tương lai}

% \subsection{Xu hướng phát sinh CTRSH tại Hội An đến năm 2030}

% Theo các số liệu lịch sử và xu hướng phát triển, lượng CTRSH tại thành phố Hội An được dự báo tiếp tục tăng trưởng trong giai đoạn 2025--2035 theo hai kịch bản. Kịch bản cơ sở (tốc độ tăng trưởng 5--7\%/năm tương ứng với tốc độ phát triển du lịch vừa phải) dự báo lượng rác đạt 130--150 tấn/ngày vào năm 2028 và 150--170 tấn/ngày vào năm 2032. Kịch bản cao (tốc độ tăng trưởng 8--10\%/năm khi Hội An mở rộng địa giới và phát triển các khu du lịch mới) dự báo lượng rác có thể vượt 180--200 tấn/ngày vào năm 2030 \cite{BaocaoHoiAn2024}.

% Với công suất thiết kế 120 tấn/ngày đêm của dự án hiện tại, nhà máy có thể xử lý đủ lượng rác đến khoảng năm 2027--2028 nếu đạt vận hành ổn định. Trong giai đoạn sau, hoặc cần mở rộng công suất bằng cách bổ sung thêm mô-đun nhiệt phân và phát điện, hoặc xây dựng nhà máy xử lý bổ sung cho phần vượt công suất. Điều này đặt ra yêu cầu thiết kế nhà máy hiện tại phải có tính dự phòng về diện tích và cơ sở hạ tầng kỹ thuật để mở rộng trong tương lai.

% Việc áp dụng mô-đun hóa trong phương án nhiệt phân -- dầu -- diesel mang lại lợi thế rõ ràng cho khả năng mở rộng. Mỗi mô-đun lò nhiệt phân (10--15 tấn/ngày) và máy phát diesel (440 kW) có thể bổ sung độc lập mà không cần dừng toàn bộ nhà máy, khác với phương án tuabin khí phải dừng và nâng cấp toàn bộ tuyến khí khi tăng công suất.

% \subsection{Tác động của phân loại rác tại nguồn đến công nghệ xử lý}

% Luật Bảo vệ môi trường 2020 yêu cầu thực hiện phân loại CTRSH tại nguồn thành ít nhất 3 loại: hữu cơ dễ phân hủy, có thể tái chế và còn lại. Lộ trình thực hiện phân loại tại nguồn tại các đô thị dự kiến hoàn thành trước năm 2025 (nội thành) và trước năm 2026 (ngoại thành và nông thôn). Khi phân loại tại nguồn được thực hiện hiệu quả, thành phần rác đầu vào nhà máy sẽ thay đổi đáng kể theo hướng thuận lợi cho công nghệ nhiệt phân.

% Cụ thể, khi phần hữu cơ dễ phân hủy được tách ra để ủ compost hoặc sản xuất khí sinh học, dòng rác còn lại cấp nhà máy sẽ có tỷ lệ nhựa, giấy và vật liệu có nhiệt trị cao tăng lên đáng kể. Độ ẩm sẽ giảm và nhiệt trị RDF sẽ tăng lên, cải thiện hiệu suất nhiệt phân và chất lượng dầu. Điều này cho thấy phương án nhiệt phân sẽ ngày càng phù hợp hơn theo thời gian khi hành vi phân loại rác của người dân được cải thiện -- một lợi thế chiến lược dài hạn so với các công nghệ xử lý rác hỗn hợp.

% \begin{table}[htbp]
% \centering
% \caption{Dự báo thành phần rác đầu vào nhà máy khi phân loại tại nguồn được cải thiện}
% \label{tab:du-bao-thanh-phan-rac}
% \setlength{\tabcolsep}{5pt}
% \resizebox{\textwidth}{!}{%
% \begin{tabular}{>{\raggedright\arraybackslash}p{0.28\textwidth}>{\centering\arraybackslash}p{0.22\textwidth}>{\centering\arraybackslash}p{0.22\textwidth}>{\raggedright\arraybackslash}p{0.20\textwidth}}
% \toprule
% \textbf{Thành phần} & \textbf{Hiện tại (không phân loại)} & \textbf{Kịch bản phân loại tốt} & \textbf{Tác động} \\
% \midrule
% Rác hữu cơ dễ phân hủy & 62,56\% & 15--20\% & Giảm mạnh \\
% Nhựa, nylon, bao bì & 10,16\% & 25--35\% & Tăng mạnh \\
% Giấy, carton & 7,42\% & 10--15\% & Tăng nhẹ \\
% Vải, gỗ, cao su & 6,48\% & 8--12\% & Tăng nhẹ \\
% Tạp chất trơ & 13,38\% & 8--12\% & Giảm nhẹ \\
% Độ ẩm trung bình & 56,3\% & 20--35\% & Giảm đáng kể \\
% Nhiệt trị RDF ước tính & 19.259 kJ/kg & 22.000--28.000 kJ/kg & Tăng đáng kể \\
% \bottomrule
% \end{tabular}%
% }
% \end{table}

% Phân tích trên cho thấy chiến lược lựa chọn công nghệ nhiệt phân phù hợp với cả điều kiện hiện tại và xu hướng dài hạn: khi phân loại tại nguồn được cải thiện, hiệu suất thu hồi năng lượng và chất lượng dầu nhiệt phân sẽ tăng theo mà không cần thay đổi cấu hình công nghệ cơ bản.

\end{document}
